\chapter{The Fourier series of the electromagnetic field emitted by a particle for the case of a circular trajectory}

Consider the case of the circular pmotion of the particle. The initial position $\mathfrak{R}_0$ is taken in the plane $Oxy$, and we denote by $\varrho_0, \varPhi_0$ the polar coordinates of $\mathfrak{R}_0$. The particle's trajectory is then given by
\begin{align}
    \mathbf{r}_0(t) = r_0 \left(\cos(\omega_L t - M\varPhi_0), \sin(\omega_L t - M\varPhi_0), 0\right)
\end{align}
(we assume the motion to be modulated by the laser field vorticity $M$ and the radius $r_0$ is the amplitude of the particle's trajectory).

The reduced velocity is given by
\begin{align}
    {\pmb\beta}(t) = \frac{d\mathbf{r}_0(t)}{dt} = r_0 \omega_L \left(-\sin(\omega_L t - M\varPhi_0), \cos(\omega_L t - M\varPhi_0), 0\right)
\end{align}

The polar angles of the unity vector ${\bf n}_0$ are $\theta_0$, $\phi_0$, i.e.
\begin{align}
    {\bf n}_0 = (\sin\theta_0\cos\phi_0, \sin\theta_0\sin\phi_0, \cos\theta_0)
\end{align}